% Generado por results_finales/r8_tiempos_adversarial.m; no editar a mano.
\begin{table}[h]
\centering
\footnotesize
\setlength{\tabcolsep}{4pt}
\caption{Prueba 2019--2025: distribución adversarial $p$ de la bola de Wasserstein frente a un mismo portafolio (el Kelly nominal con el mismo $\delta$), de modo que solo cambia la geometría. Medias sobre las 84 fechas de decisión. Masa movida: $\frac12\sum_s|p_s-1/N|$. Exceso de masa: masa adversarial en los meses de la ventana con VIX sobre su mediana (o con retorno de SPY negativo) menos la fracción de esos meses. Fechas con exceso $>0$: con umbral $10^{-8}$. Descriptivo.}\label{tab:adversarial}
\begin{tabular}{|l|r|r|r|r|}
\hline
 & \multicolumn{2}{c|}{\textbf{$c_1=(0.90,\,0.01)$}} & \multicolumn{2}{c|}{\textbf{$c_2=(0.80,\,0.01)$}} \\ \cline{2-5}
\textbf{Indicador} & \textbf{G1} & \textbf{G2} & \textbf{G1} & \textbf{G2} \\ \hline
$\rho_k$ medio & 0.0083 & 0.0083 & 0.0083 & 0.0083 \\ \hline
Masa movida (\%) & 1.10 & 1.32 & 0.97 & 1.33 \\ \hline
Número efectivo de escenarios (de 60) & 59.1 & 58.8 & 59.3 & 58.8 \\ \hline
Masa del peor escenario (\%; $1/N=1.67$) & 1.81 & 1.72 & 1.93 & 1.83 \\ \hline
Exceso de masa, VIX alto (pp) & 0.13 & $-$0.02 & 0.13 & 0.09 \\ \hline
Fechas con exceso $>0$, VIX alto (\%) & 33.3 & 6.0 & 29.8 & 11.9 \\ \hline
Exceso de masa, $r^{\mathrm{SPY}}<0$ (pp) & 0.65 & 0.39 & 0.89 & 0.42 \\ \hline
Fechas con exceso $>0$, $r^{\mathrm{SPY}}<0$ (\%) & 73.8 & 40.5 & 94.0 & 50.0 \\ \hline
\end{tabular}
\end{table}
